\usepackage[a4paper,margin=2.45cm]{geometry}
\usepackage{amsmath,amssymb,amsfonts,bm,mathtools}
\usepackage{booktabs,array,longtable,tabularx}
\usepackage{amsthm}
\usepackage{tikz}
\usepackage{pgfplots}
\pgfplotsset{compat=1.18}
\usepackage{fancyhdr}
\usepackage{placeins}
\usepackage{makeidx}
\makeindex
\usepackage{enumitem}
\usepackage{hyperref}
\usepackage{xcolor}
\usepackage[most]{tcolorbox}
\usepackage{microtype}
\usepackage{needspace}

\hypersetup{
  colorlinks=true,
  linkcolor=blue,
  urlcolor=blue,
  citecolor=blue,
  pdftitle={计算物理学简明教程：从动力学约化到现代高阶离散},
  pdfauthor={}
}

\setcounter{tocdepth}{1} % 目录显示 part、chapter 和 section
\setlength{\parindent}{2em}
\setlength{\parskip}{0.35em}

\newcommand{\dd}{\,\mathrm d}
\newcommand{\ee}{\mathrm e}
\newcommand{\ii}{\mathrm i}
\newcommand{\R}{\mathbb R}
\newcommand{\Order}{\mathcal O}
\newcommand{\Lop}{\mathcal L}
\newcommand{\Nop}{\mathcal N}
\newcommand{\Ksp}{\mathcal K}
\newcommand{\avg}[1]{\left\langle #1\right\rangle}

\newtcolorbox{chaptersummary}{
  colback=blue!3,
  colframe=blue!45!black,
  title={本章总结},
  fonttitle=\bfseries,
  boxrule=0.6pt,
  arc=2mm,
  left=2mm,
  right=2mm,
  top=1mm,
  bottom=1mm
}

\newtcolorbox{authornote}{
  colback=yellow!8,
  colframe=orange!75!black,
  title={潜在认知跳跃},
  fonttitle=\bfseries,
  boxrule=0.8pt,
  arc=2mm,
  left=2mm,
  right=2mm,
  top=1mm,
  bottom=1mm
}

\newtcolorbox{askai}{
  colback=violet!3,
  colframe=violet!55!black,
  title={你应该问 AI},
  fonttitle=\bfseries,
  boxrule=0.7pt,
  arc=2mm,
  left=2mm,
  right=2mm,
  top=1mm,
  bottom=1mm
}

\newcommand{\studyroute}[3]{%
  \par\smallskip\noindent\textbf{展开与检索。}#1\par
  \noindent\textbf{关键词：}#2。\par
  \noindent\textbf{AI 展开：}#3\par
}


\raggedbottom
\setlist{nosep,topsep=3pt}
\pagestyle{fancy}
\fancyhf{}
\fancyhead[LE,RO]{\small\nouppercase{\leftmark}}
\fancyfoot[C]{\thepage}
\setlength{\headheight}{15pt}
\newcommand{\learninggoals}[1]{\paragraph{本章学习目标}#1}
\newcommand{\conceptbridge}[3]{%
\begin{authornote}
\textbf{潜在疑问：}#1\par
\textbf{作者补充：}#2\par
\textbf{核心概念联系：}#3
\end{authornote}}
\newcommand{\aiexercise}[2]{%
\begin{askai}
\textbf{AI 展开问题：}#1\par\smallskip
\textbf{进一步展开：}#2
\end{askai}}
\newcommand{\norm}[1]{\left\lVert#1\right\rVert}
\newcommand{\jump}[1]{\left[#1\right]}
\newcommand{\avgface}[1]{\left\{#1\right\}}
\newcommand{\Poly}{\mathbb P}
\newcommand{\tr}{\operatorname{tr}}
\newcommand{\diag}{\operatorname{diag}}
\numberwithin{equation}{chapter}
\newtheorem{theorem}{定理}[chapter]
\newtheorem{proposition}[theorem]{命题}
\theoremstyle{definition}
\newtheorem{definition}[theorem]{定义}
